\subsection{The Splitting of Artinian Algebras}

In this subsection, K is a commutative ring and A a K-algebra. Let r be the radical of A. We denote the quotient algebra \(A/r\) by \(\overline{A}\) and the canonical mapping from A to \(\overline{A}\) by \(\pi\) . We are interested in the subalgebras S of A such that \(A = S \oplus r\) .

We denote by \(\Sigma\) the set of K-linear sections \(s\) of \(\pi\) satisfying \(s(\alpha\beta) = s(\alpha)s(\beta)\) for \(\alpha, \beta\) in \(\overline{\mathrm{A}}\) . Note that such a section necessarily satisfies \(s(1) = 1\) (in other words, \(s\) is a ring homomorphism): we indeed have \(s(1)^2 = s(1)\) , and \(s(1)\) is invertible because it belongs to \(1 + \mathfrak{r}\) (VIII, p.~156, Theorem 1). If \(s\) is an element of \(\Sigma\) , then the image S of \(s\) is a subalgebra of A, and we have \(\mathrm{A} = \mathrm{S} \oplus \mathfrak{r}\) . Conversely, if S is a subalgebra of A such that \(\mathrm{A} = \mathrm{S} \oplus \mathfrak{r}\) , then the restriction of \(\pi\) to S is bijective, and the inverse bijection defines an element of \(\Sigma\) with image S.

By Jacobson's theorem (loc. cit.), every element of \(1 + \mathfrak{r}\) is invertible in A. We call an inner automorphism of A of the form \(a \mapsto xax^{-1}\) with \(x \in 1 + \mathfrak{r}\) a special automorphism.

PROPOSITION 7. --- Suppose that the \((\overline{\mathrm{A}}\otimes_{\mathrm{K}}\overline{\mathrm{A}}^{\mathrm{o}})\) -module \(\overline{\mathrm{A}}\) is projective.

\begin{enumerate}
\def\labelenumi{\alph{enumi})}
\item
  Let \(S_1\) and \(S_2\) be subalgebras of A satisfying \(A = S_1 \oplus \mathfrak{r} = S_2 \oplus \mathfrak{r}\) . There exists a special automorphism of A transforming \(S_1\) into \(S_2\) .
\item
  Suppose that \(\pi\) has a K-linear section and that the radical r of A is nilpotent. Then there exists a subalgebra S of A satisfying \(A = S \oplus r\) .
\end{enumerate}

Let \(S_{1}\) and \(S_{2}\) be as in a). Let \(s_{1}\) and \(s_{2}\) be the elements of the set \(\Sigma\) corresponding to the subalgebras \(S_{1}\) and \(S_{2}\) . Let \(\varepsilon\) be the K-linear mapping from \(A \otimes_{K} A\) to \(A\) given by \(\varepsilon(a \otimes b) = ab\) . By Proposition 5 of VIII, p.~236 and Remark 1 of VIII, p.~237, there exists an element \(e = \sum_{i=1}^{r} \alpha_{i} \otimes \alpha_{i}'\) of \(\overline{A} \otimes_{K} \overline{A}\) satisfying \(\sum_{i=1}^{r} \alpha_{i} \alpha_{i}' = 1\) and \(\sum_{i=1}^{r} \alpha \alpha_{i} \otimes \alpha_{i}' = \sum_{i=1}^{r} \alpha_{i} \otimes \alpha_{i}' \alpha\) for every \(\alpha \in \overline{A}\) . Set \(x = \sum_{i=1}^{r} s_{1}(\alpha_{i}) s_{2}(\alpha_{i}')\) . We have \(\pi(x) = \sum_{i=1}^{r} \alpha_{i} \alpha_{i}' = 1\) and therefore \(x \in 1 + r\) . Let \(\alpha\) be an element of \(\overline{A}\) . We have

\[
\begin{array}{l} s _ {1} (\alpha) x = \sum_ {i = 1} ^ {r} s _ {1} (\alpha \alpha_ {i}) s _ {2} (\alpha_ {i} ^ {\prime}) = (\varepsilon \circ (s _ {1} \otimes s _ {2})) \biggl (\sum_ {i = 1} ^ {r} \alpha \alpha_ {i} \otimes \alpha_ {i} ^ {\prime} \biggr) \\ = (\varepsilon \circ (s _ {1} \otimes s _ {2})) \biggl (\sum_ {i = 1} ^ {r} \alpha_ {i} \otimes \alpha_ {i} ^ {\prime} \alpha \biggr) = \sum_ {i = 1} ^ {r} s _ {1} (\alpha_ {i}) s _ {2} (\alpha_ {i} ^ {\prime} \alpha) = x s _ {2} (\alpha). \end{array}
\]

The equality \(x^{-1}S_{1}x = S_{2}\) follows, giving assertion a).

Under the assumptions of b), suppose furthermore that \(r^{2}=0\) . In this case, the (A, A)-bimodule r is annihilated by r, and we therefore view it as an \((\overline{A},\overline{A})\) -bimodule. Choose a K-linear section \(\sigma\) of \(\pi\) . We have

\[
\alpha x = \sigma (\alpha) x \quad \text { and } \quad x \alpha = x \sigma (\alpha)\tag{19}
\]

for \(\alpha \in \overline{\mathrm{A}}\) and \(x\in \mathfrak{r}\) . Set

\[
\varphi (\alpha , \beta) = \sigma (\alpha \beta) - \sigma (\alpha) \sigma (\beta)\tag{20}
\]

for \(\alpha, \beta \in \overline{\mathbf{A}}\) . We have the relation \(\pi(\varphi(\alpha, \beta)) = \alpha\beta - \alpha\beta = 0\) for \(\alpha, \beta \in \overline{\mathbf{A}}\) . Therefore, \(\varphi\) defines an element of \(\mathrm{C}^2(\overline{\mathbf{A}}, \mathfrak{r})\) . Let \(\alpha, \beta, \gamma\) be elements of \(\overline{\mathbf{A}}\) ; in view of (19), we have

\[
\begin{array}{r l} \partial^ {2} \varphi (\alpha , \beta , \gamma) & = \alpha \varphi (\beta , \gamma) - \varphi (\alpha \beta , \gamma) + \varphi (\alpha , \beta \gamma) - \varphi (\alpha , \beta) \gamma \\ & = \sigma (\alpha) \varphi (\beta , \gamma) - \varphi (\alpha \beta , \gamma) + \varphi (\alpha , \beta \gamma) - \varphi (\alpha , \beta) \sigma (\gamma) \\ & = \sigma (\alpha) (\sigma (\beta \gamma) - \sigma (\beta) \sigma (\gamma)) - \sigma (\alpha \beta \gamma) + \sigma (\alpha \beta) \sigma (\gamma) + \sigma (\alpha \beta \gamma) \\ & - \sigma (\alpha) \sigma (\beta \gamma) - (\sigma (\alpha \beta) - \sigma (\alpha) \sigma (\beta)) \sigma (\gamma) \\ & = 0. \end{array}
\]

By Theorem 3 of VIII, p.~242, the K-module \(\mathrm{H}^{2}(\overline{\mathrm{A}},\mathfrak{r})\) is reduced to zero. Hence there exists an element \(\psi\) of \(\mathrm{C}^{1}(\overline{\mathrm{A}},\mathfrak{r})\) such that \(\partial^{1}\psi = \varphi\) , in other words, such that we have

\[
\varphi (\alpha , \beta) = \alpha \psi (\beta) - \psi (\alpha \beta) + \psi (\alpha) \beta \quad \text { for } \alpha , \beta \text { in } \overline {{\mathrm{A}}}.\tag{21}
\]

We have \(\psi(\alpha)\psi(\beta)=0\) because \(r^{2}\) is zero; from (19) and (20), we then deduce

\[
(\sigma + \psi) (\alpha \beta) = (\sigma + \psi) (\alpha) (\sigma + \psi) (\beta),\tag{22}
\]

so that the K-linear section \(\sigma +\psi\) of \(\pi\) belongs to \(\Sigma\) . Its image is a subalgebra \(S\) of \(A\) such that \(A = S + \mathfrak{r}\) .

Let us now prove the existence of S in the general case. We reason by induction on the least integer \(p \geqslant 1\) such that \(r^{p} = 0\) ; the case p = 1 is trivial. Suppose \(p \geqslant 2\) , and set \(A' = A/r^{p-1}\) . The radical \(r'\) of \(A'\) is equal to \(r/r^{p-1}\) (Proposition 5 of VIII, p.~155), so satisfies \(r'^{p-1} = 0\) , and the algebra \(A'/r'\) is isomorphic to \(\overline{A} = A/r\) and is therefore absolutely semisimple. By the induction hypothesis, there exists a subalgebra \(S'\) of \(A'\) such that \(A' = S' \oplus r'\) . Then \(S'\) is of the form \(A''/r^{p-1}\) , where \(A''\) is a subalgebra of A containing \(r^{p-1}\) , and we have

\[
\mathrm{A} = \mathrm{A} ^ {\prime \prime} + \mathfrak {r}, \quad \mathfrak {r} ^ {p - 1} = \mathrm{A} ^ {\prime \prime} \cap \mathfrak {r}.\tag{23}
\]

The algebra \(A^{\prime\prime}/r^{p-1}\) is isomorphic to \(A^{\prime}/r^{\prime}\) ; we have \((\mathfrak{r}^{p-1})^{2}=0\) , so \(r^{p-1}\) is the radical of \(A^{\prime\prime}\) . By the case we just treated, there exists a subalgebra S of \(A^{\prime\prime}\) such that \(A^{\prime\prime}=S\oplus r^{p-1}\) ; we deduce the relation \(A=S\oplus r\) from (23).

COROLLARY 1 (Wedderburn's theorem). --- Let K be a commutative field, A a K-algebra, and r the radical of A. Suppose that the K-algebra A/r is absolutely semisimple.

\begin{enumerate}
\def\labelenumi{\alph{enumi})}
\item
  Let \(S_1\) and \(S_2\) be subalgebras of A satisfying \(A = S_1 \oplus \mathfrak{r} = S_2 \oplus \mathfrak{r}\) . There exists a special automorphism of A transforming \(S_1\) into \(S_2\) .
\item
  If \(\mathfrak{r}\) is nilpotent, then there exists a subalgebra \(S\) of \(A\) satisfying \(A = S \oplus \mathfrak{r}\) .
\end{enumerate}

This follows from Proposition 7 and Theorem 2 of VIII, p.~238.

COROLLARY 2. --- Let A be a commutative algebra of finite degree over a perfect field K, and let r be its radical. There exists a unique subalgebra S of A such that \(A = S \oplus r\) . Moreover, S is isomorphic to a product of finitely many extensions of K of finite degree.

The K-algebra A/r is semisimple (VIII, p.~173, Proposition 1) and has finite degree; it is absolutely semisimple because the field K is perfect (VIII, p.~232, Theorem 1). Since the ideal r is nilpotent and A is commutative, the existence and uniqueness of S then follow from Corollary 1. Since S is semisimple and commutative and has finite degree, the last assertion is a consequence of Proposition 3 of VIII, p.~137.

Remarks. --- 1) The assumption that A/r is absolutely semisimple is essential in Corollary 1 (VIII, p.~246, Exercise 4).

\begin{enumerate}
\def\labelenumi{\arabic{enumi})}
\setcounter{enumi}{1}
\tightlist
\item
  Suppose that A is an Artinian algebra over the field K. If A is commutative, then we can show (VIII, p.~180, Exercise 9) that A is isomorphic to a product of algebras \(A_{1} \times \cdots \times A_{n}\) such that \(\mathrm{A}_{i}/\Re(\mathrm{A}_{i})\) is a field for every i. Conversely, if A is not commutative, then A may not be isomorphic to a product of algebras \(A_{1} \times \cdots \times A_{n}\) such that \(\mathrm{A}_{i}/\Re(\mathrm{A}_{i})\) is a simple ring for every i (VIII, p.~247, Exercise 5).
\end{enumerate}

